\documentclass[12pt,a4paper]{article}

\title{Documentacao Tecnica Abrangente\\Modelo de Site de Laboratorio UFOP}
\author{Projeto Laboratorio Vinculado a Departamento}
\date{\today}

\begin{document}
\maketitle
\tableofcontents
\newpage

\section{Resumo Executivo}
Este documento consolida a arquitetura, as decisoes tecnicas e os modulos implementados para o portal de laboratorio vinculado a departamento da UFOP.

A implementacao foi orientada por tres objetivos centrais:
\begin{itemize}
\item simplicidade operacional;
\item baixo consumo de recursos de servidor;
\item facilidade de manutencao por equipe institucional.
\end{itemize}

\section{Contexto e Restricoes da UFOP}
\subsection{Premissas de Infraestrutura}
O projeto foi desenhado considerando ambiente com capacidade computacional limitada, comum em servidores institucionais compartilhados.

Principais riscos considerados:
\begin{itemize}
\item CPU e memoria concorridas com outros servicos;
\item necessidade de alta disponibilidade com baixa complexidade de operacao;
\item equipe de manutencao com rotatividade e tempo reduzido para suporte.
\end{itemize}

\subsection{Diretriz de Solucao}
A diretriz adotada foi reduzir dependencias externas e componentes de runtime, priorizando:
\begin{itemize}
\item renderizacao server-side em PHP;
\item banco MySQL com estrutura simples;
\item frontend com Bootstrap e CSS local;
\item ausencia de build pipeline pesado para producao.
\end{itemize}

Essa escolha melhora previsibilidade de desempenho e simplifica suporte tecnico.

\section{Como o Laravel foi aplicado no projeto}
\subsection{Posicionamento tecnico real}
O repositorio esta no contexto de um workspace Laravel, e o projeto aproveita principios e organizacao inspirados no ecossistema Laravel, mas a entrega final para o portal do laboratorio foi otimizada para um \textbf{runtime PHP modular enxuto}, sem acoplar o ciclo completo de framework em producao.

\subsection{Elementos inspirados no modelo Laravel}
Os seguintes aspectos seguem filosofia tipica de projetos Laravel:
\begin{itemize}
\item separacao de responsabilidades entre interface, regras de negocio e persistencia;
\item configuracao centralizada e reutilizavel em \texttt{app/includes/config.php};
\item padronizacao de rotas por contexto funcional (admin, laboratorio, noticias, contato);
\item tratamento de autenticao, autorizacao e auditoria em camada central;
\item uso de seeds e inicializacao de dados para ambiente reproduzivel.
\end{itemize}

\subsection{Motivo da nao adocao full-stack Laravel no runtime final}
Por decisao arquitetural ligada a restricao de recursos, foi priorizado um arranjo com menor overhead operacional:
\begin{itemize}
\item menos servicos auxiliares e menos consumo de memoria;
\item menor custo de deploy e troubleshooting;
\item menor risco de indisponibilidade por cadeia extensa de dependencias.
\end{itemize}

Em termos praticos, a abordagem escolhida entrega os beneficios de organizacao e governanca sem impor custo de execucao elevado.

\section{Arquitetura Escolhida}
\subsection{Visao em Camadas}
A solucao segue arquitetura modular em tres camadas.

\subsubsection{Camada de Apresentacao}
\begin{itemize}
\item arquivos publicos em \texttt{app/*.php} e subpastas;
\item header e footer comuns em \texttt{app/includes/header.php} e \texttt{app/includes/footer.php};
\item tema central em \texttt{app/assets/css/theme.css}.
\end{itemize}

\subsubsection{Camada de Aplicacao}
\begin{itemize}
\item regras e servicos em \texttt{app/includes/config.php};
\item funcoes para conteudo dinamico, tema, autenticacao e permissoes;
\item funcoes para renderizacao de menu administrativo e controle de acesso.
\end{itemize}

\subsubsection{Camada de Dados}
\begin{itemize}
\item MySQL como persistencia principal;
\item settings em \texttt{site\_settings};
\item conteudos em tabelas especializadas (noticias, editais, itens dinamicos);
\item inicializacao em \texttt{mysql-init/02-seed-content.sql}.
\end{itemize}

\subsection{Fluxo de Requisicao}
Fluxo padrao de uma pagina dinamica:
\begin{enumerate}
\item request HTTP chega ao endpoint PHP;
\item endpoint chama funcoes de dominio no \texttt{config.php};
\item dados sao buscados no MySQL;
\item header e footer padronizados sao aplicados;
\item HTML final e entregue pronto ao navegador.
\end{enumerate}

\subsection{Beneficios desta arquitetura}
\begin{itemize}
\item baixo custo de execucao;
\item onboarding rapido da equipe;
\item manutencao centralizada;
\item evolucao incremental sem reescrita total.
\end{itemize}

\section{Modulos Implementados}
\subsection{Site Publico}
\begin{itemize}
\item Home com foco em noticias: \texttt{app/index.php};
\item menu principal do laboratorio;
\item paginas do laboratorio: sobre, equipe, projetos, publicacoes, cursos, parceiros, tutoriais, blog, eventos;
\item contato dinamico: \texttt{/contato/index.php};
\item noticias e editais com listagem e detalhe.
\end{itemize}

\subsection{Painel Administrativo}
\begin{itemize}
\item dashboard: \texttt{/admin/dashboard.php};
\item gestao de noticias e editais: \texttt{/admin/content.php};
\item gestao das paginas do laboratorio: \texttt{/admin/laboratorio-paginas.php};
\item gestao da pagina sobre: \texttt{/admin/laboratorio-sobre.php};
\item gestao da pagina de contato: \texttt{/admin/laboratorio-contato.php};
\item tema e paletas: \texttt{/admin/tema.php};
\item menu e logos: \texttt{/admin/menu.php};
\item usuarios e permissoes: \texttt{/admin/users.php}.
\end{itemize}

\subsection{Conteudo Dinamico com CRUD}
Implementado CRUD para paginas de laboratorio com suporte a:
\begin{itemize}
\item criacao de itens;
\item edicao;
\item remocao;
\item exibicao publica com paginacao e filtro por ano.
\end{itemize}

\subsection{Identidade Visual e Temas}
\begin{itemize}
\item motor de paletas em \texttt{site\_color\_palettes()};
\item aplicacao via variaveis CSS geradas no backend;
\item paleta \textbf{UFOP Oficial} priorizada;
\item contraste separado entre topbar, menu principal, hero e footer.
\end{itemize}

\subsubsection{UFOP Oficial}
Configuracao principal atual:
\begin{itemize}
\item cinza de apoio: \texttt{\#919191};
\item vermelho principal: \texttt{\#962038};
\item menu principal: \texttt{\#962038};
\item fundo do hero/carousel: \texttt{\#919191};
\item footer e menu do footer: \texttt{\#962038} com degrad? alinhado ao cabecalho.
\end{itemize}

\section{Seguranca, Permissoes e Auditoria}
\subsection{Autenticacao e Controle de Acesso}
\begin{itemize}
\item login administrativo com hash de senha;
\item perfis de acesso com permissoes granulares;
\item protecao de rotas administrativas por \texttt{require\_admin\_permission()}.
\end{itemize}

\subsection{Protecoes de Formulario}
\begin{itemize}
\item uso de token CSRF nos formularios administrativos;
\item validacao de entradas para operacoes de persistencia;
\item sanitizacao de saida na renderizacao HTML.
\end{itemize}

\subsection{Auditoria}
Acoes administrativas relevantes registradas em trilha de auditoria (criacao, edicao, exclusao, alteracao de tema, alteracao de usuarios).

\section{Desempenho e Operacao}
\subsection{Decisoes para desempenho em servidor limitado}
\begin{itemize}
\item renderizacao direta sem camada JS pesada;
\item minimizacao de bibliotecas externas;
\item consultas simples e previsiveis;
\item paginas com paginacao para evitar cargas grandes;
\item estrutura stateless para facilitar reinicio de container.
\end{itemize}

\subsection{Como isso melhora o funcionamento}
\begin{itemize}
\item menor tempo de resposta medio;
\item menor consumo de RAM;
\item menor risco de gargalos por dependencia externa;
\item manutencao mais simples em equipe reduzida.
\end{itemize}

\section{Infraestrutura e Deploy}
\subsection{Ambiente local e homologacao}
\begin{verbatim}
docker compose up -d --build app
\end{verbatim}

\subsection{Endpoints padrao}
\begin{itemize}
\item site: \texttt{http://localhost:8092/}
\item admin: \texttt{http://localhost:8092/admin/login.php}
\end{itemize}

\subsection{Componentes de infraestrutura}
\begin{itemize}
\item container da aplicacao PHP/Apache;
\item container MySQL;
\item volume/persistencia para dados e assets publicados.
\end{itemize}

\section{Mapa de Arquivos e Responsabilidades}
\begin{itemize}
\item \texttt{app/includes/config.php}: nucleo de negocio e acesso a dados;
\item \texttt{app/includes/header.php}: topo, logos e menu principal;
\item \texttt{app/includes/footer.php}: menu footer, calendario, botao topo;
\item \texttt{app/assets/css/theme.css}: comportamento visual global;
\item \texttt{app/admin/*.php}: administracao de conteudo e configuracoes;
\item \texttt{app/laboratorio/*.php}: paginas publicas do laboratorio;
\item \texttt{app/noticias/*.php}: noticias, editais e visualizacao;
\item \texttt{mysql-init/02-seed-content.sql}: seed inicial de conteudo.
\end{itemize}

\section{Conclusao Arquitetural}
A arquitetura adotada equilibra governanca, desempenho e simplicidade operacional. O projeto aplica principios de organizacao tipicos do ecossistema Laravel, mas com runtime PHP modular enxuto, adequado ao contexto de infraestrutura com recursos limitados da UFOP.

Esse desenho reduz custo de operacao, aumenta previsibilidade e permite evolucao incremental do portal sem comprometer estabilidade.

\end{document}
